\section{Unified restrictions on the four quadratic cases}
\label{sec:quadratic-frontier-retry}
Throughout this section $d=4+3\sqrt2$ and $D=144+96\sqrt2$.
The four quadratic rings have dimension type $(1^4,2,d^4)$, and
$N_{\mathbb Q(\sqrt2)/\mathbb Q}(D)=2304=2^8 3^2$.
All spherical and positive-center arguments use the already audited
sphericalization and Galois-existence reduction; no claim is made that
every original category is unitary or pseudo-unitary.

\subsection{The integral part of the center is pointed}
\begin{proposition}
For a complex categorification of any of the four quadratic rings,
the center contains no nontrivial Tannakian subcategory, and every
integral-dimensional central simple is invertible. The integral center
is Vec or sVec for the three commutative rings, and Vec for $R9\_2\_1975$.
\end{proposition}
\begin{proof}
The established TY module reduction gives the nonalternating bicharacter
on $V_4$. In a commutative source, let $h,h'$ stabilize its two distinct
large-object orbits. Their product $g$ is isotropic and acts freely.
The adjoint subcategory is the full source, so the universal grading
is trivial and forgetting central invertibles is injective.
The TY central-invertible formulas give self-braiding $+1$ for a lift
of $g$, and $-1$ for a lift of $h$ or $h'$. Thus a lift of $g$ would
supply a central Rep$(C_2)$.

Central de-equivariantization would have four invertibles and two
self-dual objects $X,Y$ of dimension $d$. Its rules are
\begin{gather*}
 hX=X,\quad hY=Y,\quad uX=Y,\quad uY=X,\\
 X^2=Y^2=1+h+4X+4Y,\qquad XY=u+hu+4X+4Y.
\end{gather*}
The alternatives $u^2=1$ and $u^2=h$ are exactly the already excluded
project tensors $R6\_4\_174$ and $R6\_4\_6$. All $216$ coefficients
of each identification agree. Their general exclusions are inherited
lower-rank theorems, not deductions from failure to find a model.
Hence $g$ has no central lift. There is at most one fermionic lift of
$h,h'$, since two would lift $g$. This proves the pointed-center bound.

For the noncommutative source, the corresponding hypothetical central
quotients are the excluded $R6\_4\_1$ or $R6\_4\_155$ with a pointed
permutation that cannot lift to a tensor autoequivalence. The latter
obstruction follows from the distinct uniquely lifted pointed twists
of $R6\_4\_155$. The exact source-to-quotient maps and both alternatives
are checked in the proof record. The pointed center is therefore Vec.

If Rep$(H)$ were Tannakian in the center, FP-dimension divisibility
and the norm of $D$ would force the primes of $|H|$ to lie in $\{2,3\}$.
Burnside makes $H$ solvable. A nontrivial solvable group has a nonunit
one-dimensional representation, which would give a central boson,
contradicting the preceding bounds.

Each central dimension is $a+bd$ with nonnegative integral $a,b$.
Thus its integral-dimensional simples are exactly those forgetting into
the integral TY subcategory. They form an integral braided fusion
subcategory whose dimension has prime support in $\{2,3\}$. The
categorical Burnside theorem makes it solvable, hence weakly
group-theoretical. With no Tannakian subcategory it is weakly
anisotropic, so Natale's integral weakly-anisotropic theorem makes it
pointed. This applies to the integral subcategory, not to the ambient
nonintegral center.
\end{proof}

The categorical inputs are the TY formulas of Gelaki--Naidu--Nikshych,
Proposition 4.1~\cite{RetryGNN}; ENO FP-dimension divisibility; Etingof--Nikshych--Ostrik,
\emph{Weakly group-theoretical and solvable fusion categories},
Theorem 1.6~\cite{RetryENO}; and Natale, \emph{The core of a weakly group-theoretical
braided fusion category}, Theorem 4.1~\cite{RetryNatale}. Their hypotheses and the exact
lower-rank dependencies are recorded in the linked proof package.
In particular no Morita representative has a nontrivial faithful
finite-group grading, since such a grading would give central Rep$(H)$.

\subsection{TY saturation and its sign}
Let $Y$ be the dimension-two TY object and $\epsilon$ the sign of its
parameter in a positive spherical representative. Shimizu's formula~\cite{RetryShimizu} gives
\[
 (\nu_1(Y),\ldots,\nu_8(Y))=(0,\epsilon,0,0,0,\epsilon,0,2).
\]
The center induction formula, $\nu_8(Y)=\FPdim(Y)$ and positivity force
all constituents of $I(Y)$ to have eighth-root twists. Fourier inversion
then fixes the weighted fibers. Each primitive-eighth fiber has weight
$D/4$; its only possible dimension pairs are
\[
 (2+d,2+7d)\quad\hbox{or}\quad(2+4d,2+4d).
\]
Combining this complete support condition with the integral-center
obstruction and character/induction equations gives exact separating
certificates for $\epsilon=+1$ in $R1607,R1975,R1976$. For R137 these support conditions alone do not decide the sign; Theorem~\ref{thm:r137-positive-ty-exclusion} excludes its positive sign using additional integral Verlinde restrictions.
An exact witness is evidence about those systems, not a category.

\subsection{A general sign-twisted module-functor criterion}
\begin{proposition}
Let $\mathcal D$ be a fusion subcategory, $\mathcal M$ an indecomposable
left $\mathcal D$-component in a fusion category, and $X$ a simple in it.
Suppose $h$ is an order-two invertible in $\mathcal D$,
$\underline{\operatorname{End}}_{\mathcal D}(X)=1+h$, $Xh\cong X$, and
right tensoring by $h$ moves another simple of $\mathcal M$.
Then the normalized left and right stabilizer maps on $X$ anticommute.
\end{proposition}
\begin{proof}
The actual internal end $B=1+h$ has its subgroup-algebra structure and
identifies $\mathcal M$ with right $B$-modules in $\mathcal D$, taking
$X$ to $B$. The functor $(-)h$ is a left $\mathcal D$-module functor.
The module-functor correspondence~\cite{RetryEGNO} represents it by a $B$-bimodule $P$. Since $Xh\cong X$,
its underlying right $B$-module is $B$. Its left action is an algebra
automorphism of $B$: the unit is fixed and the $h$-summand maps to
$h$ or $-h$. These exhaust the possibilities because
$\operatorname{Hom}_B(hB,B)$ is one-dimensional and $h^2=1$.
The positive action gives the regular bimodule and hence the identity
module functor, impossible because another simple is moved. The action
is therefore the negative one. Transporting it through $r:Xh\to X$
gives, with $l:hX\to X$,
\[
 r(l\otimes1_h)=-l(1_h\otimes r),
\]
with ambient associators understood. This is the required categorical
anticommutator. It uses no central lift of the based-central $h$.
\end{proof}

For each of the three commutative sources, the two pointed orbits have
different stabilizers, so right $h$ fixes the first and swaps the second.
The criterion applies to every large simple. Exact source tensors give
$\dim\operatorname{Hom}(1,X^3)=2$ in all twelve applications. The
Clifford rotation theorem then forces
\begin{equation}\label{eq:source-cubic-clifford}
 \nu_3(X)\in\{-1,-\zeta_3,-\zeta_3^2\}.
\end{equation}
The normalized relative trace to $\mathbb Q(\sqrt2)$ is consequently
$-1$ or $1/2$. No reality assertion is inferred from self-duality.
The noncommutative source has different left and right stabilizers;
this source-local criterion is not applied to it.

\subsection{The all-cochain Morita restriction}
The rank-twelve Morita reductions have pointed group $D_8$ and four
objects of dimension $d$. There are three types in their union:
commutative sources allow duality shift zero or a reflection coset;
$R1975$ allows reflection or rotation. Only the reflection type overlaps.
The earlier nonidentity quotient action and extreme $R1975$ branch
are excluded by their separate associator and pointed-automorphism
obstructions.

The forced cyclic restrictions determine the pointed $D_8$ associator
class: it restricts to class two on the rotation $C_4$, trivially on
the even $V_4$, and to the TY sign on the odd reflection groups.
The cohomology calculation of Goff--Mason--Ng~\cite[Lemmas 4.4--4.5]{RetryGMN}
shows that these cyclic restrictions detect the full class.

For a large simple, the pointed-bimodule stabilizer is
$L=D_8\times_{V_4}D_8\cong D_8\times C_2$. An explicit stabilizer
cochain, checked against every coboundary equation, has mixed pairing
$-1$ on its pure left and right central involutions $h_L,h_R$.
This is independent of the chosen solution. Indeed any other solution
differs by a $2$-cocycle $\gamma$. Its commutator with the central
$h_R$ is a character of $L$, hence is trivial on $[L,L]$; but
$h_Lh_R\in[L,L]$ and its pairing with $h_R$ equals that of $h_L$.
Thus changing the cochain cannot change the sign.

The nondegenerate stabilizer pairing makes this four-object pointed
bimodule invertible~\cite[Proposition 2.7]{RetryVV}. The cubic Clifford theorem gives
\eqref{eq:source-cubic-clifford} for each of its four large objects.
For pseudo-unitary Morita partners the weighted total indicator
\[
 v_3(\mathcal C)=\sum_T\FPdim(T)\nu_3(T)
 =\frac1D\sum_{Z\in\operatorname{Irr}\mathcal Z(\mathcal C)}
                 \FPdim(Z)^2\theta_Z^3
\]
is Morita invariant by the center indicator formula~\cite{RetryNS}. The nonunit invertibles and TY object have third
indicator zero. Subtracting the common unit and dividing by $d$
therefore transfers the sum of four large-object indicators. Its
projected value belongs to
\[
 \{-4,-5/2,-1,1/2,2\}.
\]
This covers $R1975$ as well. Neither bimodule invertibility nor
based centrality supplies a central half-braiding on the full category.

\subsection{Cubic branch restrictions and exponent consequence}
The cubic tests alone have the following feasible projected choices in their necessary center systems; these are intermediate relaxations, not the final status of a TY sign:
\begin{center}\begin{tabular}{lcl}
\toprule Source & TY sign & Four projected third indicators\\\midrule
$R137$ & $-$ & $(-1,-1,-1,-1)$\\
$R137$ & $+$ & $(1/2,1/2,1/2,1/2)$\\
$R1607$ & $-$ & $(-1,-1,-1,-1)$\\
$R1976$ & $-$ & $(1/2,1/2,1/2,1/2)$\\\bottomrule
\end{tabular}\end{center}
The $R1975$ height-six system with the aggregate condition also has
an exact witness at projected sum $2$.

For positive-sign $R137$ and negative-sign $R1976$, all other discrete
assignments are excluded: $25$ exact separating vectors and five
verified source relabellings cover all $15$ alternatives in each case.
The relabellings fix the distinguished pointed and TY objects and are
used to relabel a hypothetical categorification; no tensor
autoequivalence lifting them is assumed. Both saved matrices were
independently reconstructed entrywise.

In these two conditional branches, every actual third indicator is
$-\zeta_3$ or $-\zeta_3^2$, and $9$ divides the FS exponent. Otherwise
cubing every central twist would remove its $3$-primary part. Since
all central dimensions lie in $\mathbb Q(\sqrt2)$, the induction formula
would put $\nu_3(X)$ in a cyclotomic field of conductor coprime to $3$,
which cannot contain a nonrational cube root of unity. This contradicts
the forced indicators.

The complete proofs, explicit cochains, bounded-system semantics and
portable exact replays are curated in
\path{evidence/unresolved_retry_20261024/README.md}.
These cubic restrictions supply necessary premises, not associators or categorifiability. The following section excludes three source rings, and Theorem~\ref{thm:r137-positive-ty-exclusion} excludes the positive TY sign of the remaining ring.

